\chapter{Opis a kontext riešeného problému}\label{ch:introduction}

Webové aplikácie sú v súčasnosti neoddeliteľnou súčasťou každodenného života a nachádzajú uplatnenie v rôznych oblastiach – od administratívnych a informačných systémov až po zábavné a vzdelávacie platformy. Ich rastúca komplexnosť a požiadavky na výkon, bezpečnosť a škálovateľnosť vedú k potrebe voliť vhodnú architektúru už v počiatočných fázach vývoja. Architektúra aplikácie pritom zásadným spôsobom ovplyvňuje nielen technické parametre systému, ale aj jeho dlhodobú udržateľnosť, rozšíriteľnosť a jednoduchosť údržby.

Cieľom tejto práce je analyzovať a implementovať webovú aplikáciu na príklade hry Dáma. Projekt sa bude realizovať v dvoch odlišných architektonických prístupoch – ako monolitická aplikácia a následne ako mikroservisná architektúra. Porovnanie týchto riešení umožní lepšie pochopiť rozdiely medzi klasickým monolitickým prístupom a modernou architektúrou mikroservisov, ako aj vyhodnotiť ich vplyv na výkon, údržbu, rozšíriteľnosť a flexibilitu softvérových systémov.

\section{Kontext riešeného problému}

Hra Dáma je klasická stolová hra pre dvoch hráčov, ktorá vyžaduje interaktívnu prácu s hernou doskou, evidenciu hráčov, zaznamenávanie ťahov a skóre. V súčasnej monolitickej implementácii sú všetky komponenty aplikácie – logika hry, správa používateľov, komentáre a skóre – sústredené v jednej aplikácii. Dáta o hráčoch a priebehu hry sú uložené v session objektoch alebo priamo v databáze, pričom celý systém funguje ako jednotná aplikácia.

Monolitický prístup má viaceré výhody – jednoduchšiu implementáciu, nižšiu zložitosť pri nasadzovaní a centrálnu správu kódu. Na druhej strane však môže spôsobovať problémy pri škálovaní, údržbe a rozširovaní funkcionality. Ak sa napríklad zvyšuje počet používateľov alebo požiadaviek, monolitická aplikácia musí byť škálovaná ako celok, čo je často neefektívne a vedie k vyšším nákladom na prevádzku aj údržbu.

Z tohto dôvodu je vhodné preskúmať mikroservisnú architektúru, v ktorej je aplikácia rozdelená na viacero nezávislých služieb. Každá služba vykonáva konkrétnu funkciu – napríklad správa používateľov, logika hry, správa skóre či komentárov. Tento prístup umožňuje lepšiu škálovateľnosť, nezávislé nasadzovanie jednotlivých komponentov a väčšiu flexibilitu pri rozvoji systému. Komunikácia medzi službami prebieha prostredníctvom REST API, ktoré zabezpečuje výmenu dát medzi backendom a frontendovou časťou aplikácie.

Na ukladanie informácií o hráčoch počas hernej relácie, ako sú ich prezývky a zvolené avatary, sa v mikroservisnej verzii používa Redis, ktorý nahrádza pôvodné riešenie založené na Session Scope. Redis poskytuje rýchly prístup k údajom v pamäti a umožňuje efektívnu správu relácií používateľov bez potreby uchovávania týchto informácií v pamäti servera.

Samotné herné pole a logika hry sú naďalej spracovávané pomocou objektov v pamäti aplikácie, podobne ako v pôvodnej monolitickej verzii. Každá pozícia na doske je reprezentovaná objektom Tile a figúrky objektom Man s príslušným stavom (TileState). Inicializácia dosky prebieha v pamäti a všetky herné operácie, ako pohyb figúrok alebo kontrola pravidiel, sa vykonávajú priamo nad týmito objektmi.

Perzistentné údaje, ako sú informácie o hráčoch, skóre, komentáre a hodnotenia, sú uchovávané v PostgreSQL, čo zabezpečuje dlhodobé a spoľahlivé uloženie údajov.

Celý systém je kontajnerizovaný pomocou Dockeru, čo zjednodušuje nasadzovanie, izoláciu a správu jednotlivých služieb v lokálnom prostredí. Teoreticky by projekt mohol byť rozšírený aj o Kubernetes, ktorý sa v praxi používa na orchestráciu kontajnerov a automatické škálovanie, avšak v rámci tejto práce je projekt navrhnutý ako lokálne riešenie určené na beh na jednom zariadení.

Projekt využíva API Gateway ako jednotný vstupný bod pre frontend (Thymeleaf) a backendové mikroservisy. Gateway umožňuje smerovanie požiadaviek, základnú autentifikáciu a aplikáciu CORS. Aj keď projekt beží lokálne na jednom zariadení, jeho použitie demonštruje praktickú implementáciu mikroservisnej architektúry a ukazuje, ako by sa systém mohol rozšíriť v produkčnom prostredí.

V kontexte mikroservisov je vhodné spomenúť aj ďalšie technológie bežne používané v moderných systémoch, ako Apache Kafka pre asynchrónnu komunikáciu a streamovanie dát, či Prometheus a Grafana pre monitoring a vizualizáciu výkonu. Tieto technológie nie sú priamo implementované v projekte, ale poskytujú širší pohľad na ekosystém mikroservisných riešení.

Tento projekt preto umožňuje praktické porovnanie dvoch architektonických prístupov a poskytuje konkrétny kontext, ako sa jednotlivé technológie správajú pri implementácii interaktívnej webovej aplikácie určenej pre lokálnu hru dvoch hráčov na jednom zariadení. Porovnanie sa zameriava nielen na výkon, ale aj na zložitosť vývoja, udržiavateľnosť a flexibilitu rozširovania systému, čo predstavuje kľúčové faktory pri rozhodovaní o architektúre modernej webovej aplikácie.